\chapter*{ตารางแผนปฏิบัติการ 13 การทดลอง และการเชื่อมโยง CLOs สู่ PLOs}
\addcontentsline{toc}{chapter}{ตารางแผนปฏิบัติการ 13 การทดลอง และการเชื่อมโยง CLOs สู่ PLOs}

ปฏิบัติการในรายวิชา 4011111 ฟิสิกส์สำหรับจุลชีววิทยา จัดการเรียนการสอนสัปดาห์ละ 3 ชั่วโมง ควบคู่กับชั่วโมงบรรยาย 2 ชั่วโมง โดยมีสัดส่วนคะแนนภาคปฏิบัติการรวม \textbf{50 คะแนน} (แบ่งเป็น รายงานปฏิบัติการกลุ่ม 20 คะแนน, ทักษะการใช้อุปกรณ์ ความปลอดภัย และความซื่อสัตย์ในการบันทึกข้อมูล 15 คะแนน และโครงงานย่อยสัปดาห์ที่ 14 อีก 15 คะแนน) ดังรายละเอียดใน\mbox{ตารางที่} 1

\begin{table}[H]
\centering
\caption{แผนการจัดปฏิบัติการ 13 การทดลอง (ตลอด 15 สัปดาห์) และการเชื่อมโยงกับผลลัพธ์การเรียนรู้รายวิชา (CLOs)}
\scriptsize\setlength{\tabcolsep}{3.5pt}
\begin{tabularx}{\textwidth}{c c >{\raggedright\arraybackslash}p{4.4cm} >{\raggedright\arraybackslash}X c}
\toprule
\textbf{สัปดาห์} & \textbf{รหัสแล็บ} & \textbf{ชื่อปฏิบัติการทดลอง (3 ชม./สัปดาห์)} & \textbf{เครื่องมือจุลชีววิทยาและทักษะหลักที่ฝึก} & \textbf{CLOs ที่วัด} \\
\midrule
1 & ปฏิบัติการที่ 1 & การวัดละเอียดและทักษะการสอบเทียบไมโครปิเปตและเครื่องชั่งดิจิทัล & เวอร์เนียร์ ไมโครมิเตอร์ เครื่องชั่ง 4 ตำแหน่ง และ Micropipette (P200, P1000) & CLO1, CLO3, CLO4 \\
\addlinespace[1.5pt]
2 & ปฏิบัติการที่ 2 & การเคลื่อนที่ ความเร็ว และการสังเกตการแพร่และ Brownian Motion & การเขียนกราฟความเร็ว และการแยกแยะ Brownian Motion จากการว่ายจริงของเซลล์ & CLO1, CLO3, CLO4 \\
\addlinespace[1.5pt]
3 & ปฏิบัติการที่ 3 & กฎของนิวตันและการวัดความหนืดของของเหลวชีวภาพตามกฎของสโตกส์ & การวัดอัตราเร็วสุดท้ายของการตกตะกอน ($v_t$) และความหนืดสารละลายกลีเซอรีน & CLO1, CLO2, CLO3 \\
\addlinespace[1.5pt]
4 & ปฏิบัติการที่ 4 & การอนุรักษ์พลังงานกลและการคำนวณค่า RCF ของเครื่องปั่นเหวี่ยง & เครื่อง Centrifuge การแปลงค่า RPM เป็น $\times g$ และการจัดสมดุลโรเตอร์ & CLO1, CLO2, CLO4 \\
\addlinespace[1.5pt]
5 & ปฏิบัติการที่ 5 & ความจุความร้อนจำเพาะและการตรวจสอบเสถียรภาพตู้อบเลี้ยงเชื้อ & แคลอริมิเตอร์ อ่างน้ำควบคุมอุณหภูมิ (Water Bath) และตู้อบ Incubator $37\,^\circ\text{C}$ & CLO1, CLO2, CLO3 \\
\addlinespace[1.5pt]
6 & ปฏิบัติการที่ 6 & ความดันไอ จุดเดือด และวิเคราะห์รอบการทำงานหม้อนึ่งความดันไอ & หม้อนึ่งความดันไอ Autoclave ($121\,^\circ\text{C}$, $15\text{ psi}$) และ Autoclave Tape & CLO1, CLO2, CLO4 \\
\addlinespace[1.5pt]
7 & \multicolumn{3}{>{\raggedright\arraybackslash}p{9.0cm}}{\textbf{สัปดาห์ทบทวนสรุปสูตรคำนวณฉบับเข้าใจง่ายและสอบวัดผลกลางภาค (เนื้อหาสัปดาห์ที่ 1--6)}} & CLO1, CLO2 \\
\addlinespace[1.5pt]
8 & ปฏิบัติการที่ 7 & ไฟฟ้าสถิตและกลไกแรงดึงดูดทางไฟฟ้าของสีย้อมเซลล์จุลินทรีย์ & อิเล็กโทรสโคป และการเปรียบเทียบการยึดเกาะของสีย้อมประจุบวกกับประจุลบ & CLO1, CLO3, CLO4 \\
\addlinespace[1.5pt]
9 & ปฏิบัติการที่ 8 & กฎของโอห์มและแบบจำลองการแยกโมเลกุลด้วย Gel Electrophoresis & มัลติมิเตอร์ดิจิทัล แหล่งจ่ายไฟกระแสตรง และชุดรันเจลอะกาโรส (Run to Red) & CLO1, CLO2, CLO3 \\
\addlinespace[1.5pt]
10 & ปฏิบัติการที่ 9 & การสะท้อน การหักเหของแสง และการวัดค่าบริกซ์ของกระบวนการหมัก & ชุดเลเซอร์วัดมุมหักเหของสเนลล์ และเครื่อง Hand Refractometer วัดน้ำตาลหมักยีสต์ & CLO1, CLO3, CLO4 \\
\addlinespace[1.5pt]
11 & ปฏิบัติการที่ 10 & เลนส์นูน กล้องจุลทรรศน์ และบทบาทของน้ำมันช่วยส่องที่หัวเลนส์ 100x & กล้องจุลทรรศน์แบบใช้แสง การคำนวณค่า $d_{\min}$ และการใช้ Immersion Oil ($n=1.515$) & CLO1, CLO2, CLO3 \\
\addlinespace[1.5pt]
12 & ปฏิบัติการที่ 11 & การจำลองครึ่งชีวิตกัมมันตภาพรังสีและประสิทธิภาพรังสี UV-C ฆ่าเชื้อ & แบบจำลองการลดลงครึ่งชีวิต ($T_{1/2}$) และกฎกำลังสองผกผันของความเข้มรังสี UV & CLO1, CLO2, CLO4 \\
\addlinespace[1.5pt]
13 & ปฏิบัติการที่ 12 & การวัดค่าการดูดกลืนแสง $\text{OD}_{600}$ ประเมินการเจริญเติบโตของจุลินทรีย์ & เครื่อง Spectrophotometer หลอด Cuvette และการสร้างกราฟมาตรฐานความขุ่นเซลล์ & CLO1, CLO2, CLO4 \\
\addlinespace[1.5pt]
14 & ปฏิบัติการที่ 13 & โครงงานย่อย (Mini-Project) ฟิสิกส์ในเครื่องมือห้องปฏิบัติการจุลชีววิทยา & การสังเคราะห์องค์ความรู้และนำเสนอโปสเตอร์/สาธิตการใช้เครื่องมืออย่างปลอดภัย & CLO2, CLO3, CLO4 \\
\addlinespace[1.5pt]
15 & \multicolumn{3}{>{\raggedright\arraybackslash}p{9.0cm}}{\textbf{สัปดาห์สรุปองค์ความรู้รวบยอดและสอบวัดผลปลายภาค (เนื้อหาสัปดาห์ที่ 8--14)}} & CLO1, CLO2 \\
\bottomrule
\end{tabularx}
\end{table}
